\section{Co-energy is a latent surface inferred from its slopes}
\label{sec:potential}
\subsection{The experiment measures slopes, not absolute height}
Voltage reconstruction provides current locations and two flux components.
For a conservative static magnetic relation define
\begin{equation}
 \boxed{d\coenergy=\psid\,di_d+\psiq\,di_q,\qquad
        \nabla_{\current}\coenergy=\flux.}\label{eq:potential-gradient}
\end{equation}
The current plane is the horizontal domain of a scalar landscape.
$\psi_d$ is its slope along d current and $\psi_q$ its slope along q current.
The reconstructed flux vector thus provides local slope information, from
which a compatible landscape can be inferred. It does not provide an
independently measured co-energy height.

For amplitude-invariant dq conventions, use physical three-phase co-energy
divided by $3/2$, so \cref{eq:potential-gradient} has exactly the usual dq
flux components. $W'$ has units $\si{\weber\ampere}=\si{\joule}$ under
this fixed normalization. It is not interchangeable with magnetic energy;
where a local Legendre relation exists, their distinction is the usual energy
versus co-energy distinction. The fixed PM state and reference offset remain
part of the model.
\begin{figure}[tb]
\centering
\begin{tikzpicture}
\begin{axis}[publication 3d,width=.65\textwidth,height=6.2cm,
 xlabel={$x=i_d/I_*$},ylabel={$y=i_q/I_*$},zlabel={$\normalizedcoenergy=W'/(\psi_*I_*)$},
 domain=-1:1,y domain=-1:1,samples=21,samples y=21]
 \addplot3[scalar surface]
  {x+.3*x^2+.5*y^2-.01*x^4-.02*y^4-.03*x^2*y^2+.02*x*y^2};
\end{axis}
\end{tikzpicture}
\caption{Co-energy as a scalar landscape, using the baseline synthetic
polynomial. Flux measures local slopes; differential inductance measures
curvature. The PM term tilts the landscape along d current, while its q
reflection remains intact. The surface height is inferred, not measured.}
\label{fig:landscape}
\end{figure}


\subsection{Relative co-energy and gauge freedom}
For any admissible path $\mathcal C$ from $\current_0$ to $\current$,
\begin{equation}
 \boxed{\coenergy(\current)-\coenergy(\current_0)
       =\int_{\mathcal C}\flux\cdot d\current.}\label{eq:path-integral}
\end{equation}
Adding a constant changes no slope:
$\nabla(\coenergy+C)=\nabla\coenergy$. If the origin lies in the domain, choose
$\coenergy(0,0)=0$. This is a gauge, not a measurement of absolute energy.
On a connected domain with a fully known conservative field the potential
is unique up to this constant. Finite samples can leave additional unobserved
shape directions; a gauge alone does not eliminate them.

\subsection{Two paths reveal integrability}
On a rectangle connecting the origin to $(i_d,i_q)$, compare
\begin{align}
 W'_1&=\int_0^{i_d}\psi_d(\xi,0)\,d\xi
        +\int_0^{i_q}\psi_q(i_d,\eta)\,d\eta,\\
 W'_2&=\int_0^{i_q}\psi_q(0,\eta)\,d\eta
        +\int_0^{i_d}\psi_d(\xi,i_q)\,d\xi.
\end{align}
The same endpoint must have the same height. Going around a small rectangle
therefore produces no accumulated co-energy. For a $C^1$ field the local
condition is
\begin{equation}
 \boxed{\partial_{i_q}\psid=\partial_{i_d}\psiq.}\label{eq:integrability}
\end{equation}
On an open simply connected domain this condition is also sufficient for a
potential. On a domain with holes, zero circulation around noncontractible
loops is additionally required. Smooth local cross-derivative equality alone
does not settle that global topology.

Define the curl-like residual
\begin{equation}
 r_{\rm curl}=\partial_{i_d}\psiq-\partial_{i_q}\psid.
\end{equation}
Green's theorem for a counterclockwise boundary gives
\begin{equation}
 \boxed{\oint_{\partial\mathcal A}\flux\cdot d\current
        =\iint_{\mathcal A}r_{\rm curl}\,di_d\,di_q.}
 \label{eq:green}
\end{equation}
For positive $i_d,i_q$, the d-then-q path followed by the reverse alternative
path is counterclockwise, so $W'_1-W'_2$ equals this integral over their
rectangle. Circulation means no single-valued static scalar surface can
reproduce that field exactly. It does not specify the cause of inconsistency.
\begin{figure}[tb]
\centering
\begin{tikzpicture}[scale=1.05]
 \fill[derived light quantity] (0,0) rectangle (3,1.8);
 \draw[coordinate vector] (-.3,0)--(3.7,0) node[right]{$i_d$};
 \draw[coordinate vector] (0,-.3)--(0,2.4) node[above]{$i_q$};
 \draw[strong line,reference quantity,direction arrow] (0,0)--(3,0);
 \draw[strong line,reference quantity,direction arrow] (3,0)--(3,1.8);
 \draw[strong line,estimated quantity,alternative pattern,direction arrow] (0,0)--(0,1.8);
 \draw[strong line,estimated quantity,alternative pattern,direction arrow] (0,1.8)--(3,1.8);
 \node[operating point] at (0,0) {};
 \node[operating point] at (3,1.8) {};
 \node[above right] at (3,1.8) {same endpoint};
 \node[reference quantity,below] at (1.5,0) {$\mathcal C_1$: d then q};
 \node[estimated quantity,above] at (1.5,1.8) {$\mathcal C_2$: q then d};
 \node at (1.5,.9) {$\mathcal A$};
 \node[annotation,text width=66mm] at (7,.95)
  {$W'_1-W'_2=\displaystyle\iint_{\mathcal A}r_{\rm curl}\,di_d\,di_q$\\[5pt]
   Conservative field: both routes agree.\\
   Circulation: no single-valued height fits both.};
\end{tikzpicture}
\caption{Integrability as path agreement. Follow the reference quantity path and reverse
the estimated quantity path to form the counterclockwise loop used by Green's theorem.
Choosing only one measured integration path can hide inconsistency.}
\label{fig:integration}
\end{figure}


\subsection{Why path integration is explanatory rather than the primary fit}
Noise, resistance mismatch, voltage error, current bias, interpolation, or
excluded dynamic physics can make measured paths disagree. Selecting one
arbitrary path hides that disagreement and accumulates sampling error.
A coherent constant rotor-angle rotation, however, can retain integrability
while rotating the symmetry axis; the companion derives this distinction.
It must not be listed as an automatic cause of nonzero curl.
We therefore fit gradients directly, using all samples jointly, rather than
first creating noisy co-energy observations by integration.

\subsection{A linear PM example as a check}
For $\psi_d=\psi_{\rm pm}+L_di_d$ and $\psi_q=L_qi_q$, either path yields
\begin{equation}
 \coenergy=\psi_{\rm pm}i_d+\tfrac12L_di_d^2+\tfrac12L_qi_q^2+C.
\end{equation}
Its d derivative is $\psi_{\rm pm}+L_di_d$ and its q derivative is $L_qi_q$.
The PM supplies a linear d tilt, not an obstacle to a scalar representation.
Saturation generalizes the quadratic curvature while leaving the gradient
relationship intact. We next derive reflection without relying on this
linear special case.
